\begin{lstlisting}[language=mathematica]
changmoduli = ({\[Tau] -> E^-\[Phi] vol^(1/2), \[Rho] -> vol^( 1/3)} /. {vol -> (\[Tau] Exp[\[Phi]])^(3/2)});
Print["VH3="]
VH3 = (AH3 Exp[2 \[Phi]])/\[Rho]^6 /. changmoduli /. {\[Phi] -> Log[1/s]}
Print["V03="]
VO3 = AO3 PowerExpand[(\[Rho]^((q - 6)/2) \[Tau]^-3) /. changmoduli] /. {\[Phi] -> Log[1/s]} /. {q -> 3}
Print["AF5="]
VF5 = AF5 PowerExpand[(\[Rho]^(3 - p) \[Tau]^-4) /. changmoduli] /. {\[Phi] -> Log[1/s]} /. {p -> 5}
VF5 = AF5/\[Tau]^4;
Print["VF3="]
VF3 = (AF3 Exp[4 \[Phi]])/\[Rho]^6 /. changmoduli /. {\[Phi] -> Log[1/s]}
Print["VD5="]
VR6 = AR6 PowerExpand[(\[Rho]^-1 \[Tau]^-2) /. changmoduli];
VD5 = AD5/(Sqrt[s] \[Tau]^(5/2))
Vins = AA Exp[-aa \[Tau]];
Print["Vreal = "]
Vreal = VH3 + VF3 + VD5 + VF5 + VO3 + 0 Vins
var = {s, \[Tau]};
NM = Length[var];
(*
Potential descriptors section:
mij .......... is the hessian of the potential        |
detmij ....... is the determinant of the hessian      |
trmij ........ is the trace of the hessian            |
diV .......... is the divergent of the potential      |
egval ........ stores the eigenvalues of the hessian  |
pfun ......... the function which likely evaluates    |
               the potential and extract the egvals   |
               which later are used to return a value |
               $\lambda_1,2$ $\leq$ 0. This value is  |
                 the sum of both egvals or 0 for each |
                 semidefinite positive egval          |
 *)
Print["Hessian of the potential = "]
mij = Simplify[ Table[D[Vreal, var[[i]], var[[j]]], {i, 1, NM}, {j, 1, NM}]]
detmij = Det[mij];
trmij = Tr[mij];
diV = Simplify[Table[D[Vreal, var[[i]]], {i, 1, NM}]];
egval = Eigenvalues[ Table[D[Vreal, var[[i]], var[[j]]], {i, 1, 2}, {j, 1, 2}]];
pfun = Function[{s, \[Tau]}, 10^4 If[egval[[1]] < 0, egval[[1]], 0] + 10^4 If[egval[[2]] < 0, egval[[2]], 0]];
\end{lstlisting}

